%% kapa: exam=1 round=32 year=2021 period=1 form=A subject=real_estate subject_ko=부동산학원론 seq=21 number=101
%% source: Q-Net 감정평가사 기출문제 (cst00309) artlSeq=5208397 file=extracted/1008/5208397/A형/2021년도 제32회 감정평가사 1차 1교시 A형.pdf
%% conversion: kordoc 4.24.51 -> tools/gen1.py (auto); review=proofread-sonnet; flags=none
\begin{question}{101}
부동산 증권에 관한 설명으로 옳지 않은 것은?
\begin{choices}
  \item MPTS(Mortgage Pass-Through Securities)는 지분을 나타내는 증권으로서 유동화기관의 부채로 표기되지 않는다.
  \item CMO(Collateralized Mortgage Obligation)는 동일한 저당풀(mortgage pool)에서 상환 우선순위와 만기가 다른 다양한 증권을 발행할 수 있다.
  \item 부동산개발PF ABCP(Asset Backed Commercial Paper)는 부동산개발PF ABS(Asset Backed Securities)에 비해 만기가 길고, 대부분 공모로 발행된다.
  \item MPTS(Mortgage Pass-Through Securities)는 주택담보대출의 원리금이 회수되면 MPTS의 원리금으로 지급되므로 유동화기관의 자금관리 필요성이 원칙적으로 제거된다.
  \item MBB(Mortgage Backed Bond)는 주택저당대출차입자의 채무불이행이 발생하더라도 MBB에 대한 원리금을 발행자가 투자자에게 지급하여야 한다.
\end{choices}
\end{question}
\answer{3}
